
\begin{frame}{Appendix: competencies as theoretical context}
\centering
\resizebox{.74\textwidth}{!}{% Adapted from the user's updated framework. The key is on the next slide.
\begingroup
\tikzset{
original/.style={line width=1.1pt},
relation/.style={line width=1.5pt,{Triangle[length=2mm]}-{Triangle[length=2mm]}},
labelbox/.style={draw=black,line width=.4pt,fill=white,inner sep=3pt,align=left,font=\fontsize{10}{11}\selectfont},
maintext/.style={align=center,font=\bfseries\fontsize{11}{12}\selectfont,inner sep=0pt},
badge/.style={circle,draw=competencygreen,fill=white,text=competencygreen,minimum size=4.5mm,inner sep=0pt,font=\bfseries\small},
feedback/.style={draw=competencygreen,line width=1.1pt,-{Triangle[length=2mm]}}
}
\begin{tikzpicture}[x=.01cm,y=-.01cm]
% Learner, scientific model, and phenomenon.
\draw[original] (85,686)
  .. controls (60,660) and (46,626) .. (46,591)
  .. controls (46,504) and (122,433) .. (216,433)
  .. controls (310,433) and (385,504) .. (385,591)
  .. controls (385,626) and (374,659) .. (350,686)
  .. controls (373,714) and (385,749) .. (385,785)
  .. controls (385,872) and (310,942) .. (216,942)
  .. controls (122,942) and (46,872) .. (46,785)
  .. controls (46,748) and (60,713) .. (85,686);
\draw[original,densely dotted] (85,686)--(350,686);
\draw[original] (650,316)
  .. controls (637,295) and (630,272) .. (630,248)
  .. controls (630,161) and (706,90) .. (800,90)
  .. controls (894,90) and (970,161) .. (970,248)
  .. controls (970,273) and (964,296) .. (952,316)
  .. controls (963,338) and (969,360) .. (969,384)
  .. controls (969,471) and (893,541) .. (800,541)
  .. controls (706,541) and (629,471) .. (629,384)
  .. controls (629,359) and (637,335) .. (650,316);
\draw[original,densely dotted] (650,316)--(952,316);
\draw[original] (1180,432) rectangle (1518,942);
\draw[original] (1180,686)--(1518,686);

% Relations. Orange labels describe scientific processes, not assessed skills.
\draw[relation] (412,569)--(1141,569);
\draw[relation] (412,808)--(1141,808);
\draw[relation,alignmentblue] (385,571)--(632,246);
\draw[relation,alignmentblue] (385,786)--(637,466);
\draw[relation,modellingorange] (970,246)--(1180,578);
\draw[relation,modellingorange] (950,470)--(1180,812);

\node[maintext] at (216,380) {Learner\\Mental models};
\node[maintext] at (797,69) {Physics model};
\node[maintext] at (1348,407) {Worldly entity};
\node[maintext] at (217,561) {Gestalt};
\node[maintext] at (217,815) {Functionality};
\node[maintext] at (800,204) {Selected\\Forms and\\Structures};
\node[maintext] at (800,402) {Selected\\Functions};
\node[maintext] at (1348,561) {Sensory\\patterns};
\node[maintext] at (1348,815) {Functions};
\node[labelbox,text=alignmentblue] at (492,430) {Structural\\alignment};
\node[labelbox,text=alignmentblue] at (504,694) {Functional\\alignment};
\node[labelbox,text=modellingorange] at (1075,385) {Structural\\modelling};
\node[labelbox,text=modellingorange] at (1075,663) {Functional\\modelling};
\node[labelbox,align=center] at (777,605) {Fidelity of Gestalt};
\node[labelbox,align=center] at (792,855) {Functional fidelity};

% Green markers denote the learner's activities at each relation.
% 3 and 4 mean mapping model to phenomenon, not constructing scientific models.
\node[badge] at (606,410) {1};
\node[badge] at (620,675) {2};
\node[badge] at (1138,319) {3};
\node[badge] at (1138,597) {4};
\node[badge] at (946,605) {6};
\node[badge] at (960,855) {6};

% Higher-order competency 5: compact upper-left annotation and smooth return.
\draw[feedback,line cap=round,line join=round] (650,165)
  .. controls (540,-20) and (20,-20) .. (20,185)
  .. controls (20,290) and (216,285) .. (216,340);
\node[anchor=north west,inner sep=3pt,align=left,text width=5cm,
  text=competencygreen,font=\fontsize{7.5}{9.5}\selectfont] at (75,80)
  {\textbf{Model selection and transformation}\\[3pt]
   Choose and coordinate models and\\
   representations for a purpose; use them\\
   to inform one's own understanding};
\node[badge] at (294,37) {5};
% Evidence-based revision returns to both components of the mental model.
\draw[feedback] (1350,950)
  .. controls (1350,1006) and (1310,1020) .. (1230,1020)
  -- (80,1020)
  .. controls (-62,1020) and (-65,956) .. (-65,861)
  -- (-65,648)
  .. controls (-65,590) and (-28,590) .. (44,590);
\draw[feedback] (-65,850)
  .. controls (-65,793) and (-20,793) .. (43,793);
\node[fill=white,inner sep=4pt,align=center,text=competencygreen,
  font=\fontsize{8.2}{10.6}\selectfont] at (740,1020)
  {\textbf{Learner competency: updating the mental model}\\
   Test and evaluate; revise Gestalt and Functionality in response to evidence};
\node[badge] at (210,1020) {7};
\end{tikzpicture}
\endgroup
}
% \source{Green markers: learner activities. Orange relations: scientific modelling processes.}
\note{The framework distinguishes the learner's mental model, the physics model, and the worldly entity. Fidelity describes the learner's beliefs; alignment describes correspondence with expert consensus. Orange relations denote scientific modelling processes involved in developing models and scientific consensus. Green markers 3 and 4 denote the learner's identification and comparison of correspondences across those relations, not the construction of scientific models. Return arrow 5 connects model selection and transformation to the learner's understanding. Return arrow 7 connects testing and evaluation to evidence-based revision of both Gestalt and Functionality. The seven competencies are theoretical context, not measured outcomes of this study.}
\end{frame}

\begin{frame}{Appendix: examples of modeling practices}
\footnotesize\renewcommand{\arraystretch}{1.04}
\begin{tabularx}{\textwidth}{@{}P{.025\textwidth}P{.32\textwidth}Y@{}}
\toprule
& \textbf{Learner competency} & \textbf{Illustrative practice, not an added task}\\
\midrule
1 & Structural mapping:\newline mental--physics & Compare one's imagined polarization ellipse with the physics model's selected structure.\\
2 & Functional mapping:\newline mental--physics & Compare one's explanation of phase effects with the physics model's functional relationships.\\
3 & Structural modelling:\newline physics--world mapping & Map a depicted polarizer's orientation to the transmission axis of a physical polarizer.\\
4 & Functional modelling:\newline physics--world mapping & Relate the model's predicted transmission to changes in physical light intensity.\\
5 & Model selection and transformation & Use Jones notation to calculate a transformation and the sphere to interpret it geometrically.\\
6 & Metacognitive reflection & Examine beliefs about the sphere's visual realism and functional adequacy, including its abstractions and limits.\\
7 & Evidence-based revision & Use a mismatch between prediction and applet feedback to revise Gestalt, Functionality, or both.\\
\bottomrule
\end{tabularx}
% \source{Examples adapted to the updated framework. Model competencies are not assessed in this study.}
\note{Rows 1 and 2 concern correspondences between the learner's mental model and the physics model. Rows 3 and 4 concern learner mapping between the physics model and the phenomenon. Their labels follow the supplied framework; the examples do not imply that learners construct scientific consensus or that applet feedback constitutes an independent physical experiment. Row 5 includes purposeful selection, transformation, and coordination of representations to inform one's own understanding. Row 6 concerns reflection on fidelity beliefs, while expert alignment is a separate comparison with expert-selected content. Row 7 updates both possible components of the mental model in response to evidence. These examples clarify the framework; they are not additional activities or interview prompts. Think-aloud observations remain limited to spontaneous reasoning during the existing applet tasks.}
\end{frame}



